\documentclass[12pt]{article}
\usepackage[utf8]{inputenc}
\usepackage[T1]{fontenc}
\usepackage{hyperref}
\title{Slovakian Sign Language Language Resources for K-12 Curriculum Development}
\author{Compiled by Claude Sonnet 5.5 (claude-sonnet-5-5) for Andr\'as Kornai}
\date{}
\begin{document}
\maketitle

\section{Introduction}
Slovak Sign Language (slovensk\'y posunkov\'y jazyk, ISO 639-3: svk; Glottocode slov1263) is the community sign language of deaf people in Slovakia. It has legal recognition and was codified in 2020. The search found legal texts, a multilingual online dictionary and press coverage, but no specific school textbook series.

\section{Primary Educational Resources}
No sign-language textbooks or curricula were found in the search results. A 2017 amendment to the Slovak sign-language law (Act 151/2017) establishes a right to education in Slovak Sign Language and provides for teaching materials in it:
\begin{itemize}
\item Act 151/2017 (consolidated text of 1 September 2017): \url{https://static.slov-lex.sk/static/SK/ZZ/2017/151/20170901.html}
\end{itemize}

\section{Linguistic and Reference Materials}
\begin{itemize}
\item Dictio, a multilingual sign and spoken language dictionary (Masaryk University) covering Slovak Sign Language, with about 4,809 entries according to one source. Some content may require registration: \url{https://www.sign-lang.uni-hamburg.de/lr/compendium/lex/dictio.html}
\item Presenting Dictio, an online sign and spoken language dictionary: \url{https://www.phil.muni.cz/vyzkum/publikace/prehled/2223960}
\item EU Accessible Europe Centre note on Dictio: \url{https://accessible-eu-centre.ec.europa.eu/dictio-multilingual-dictionary-focused-sign-language-czech-republic_en}
\end{itemize}

\section{Religious and Community Resources}
\begin{itemize}
\item The European Union of the Deaf page mentioned the Asociácia Nepočujúcich Slovenska (ANEPS): \url{https://eud.eu/?p=2603}
\end{itemize}
The Slovak Association of the Deaf (Slovenský zväz nepočujúcich) came up in coverage, but no page of its own was found.

\section{Digital and Audio Resources}
\begin{itemize}
\item The Dictio dictionary above is the main online resource found.
\item The Spectator (SME) article ``Slovak sign language does not help every deaf person'': \url{https://spectator.sme.sk/c/22805953/slovak-sign-language-does-not-help-every-deaf-person.html}
\end{itemize}

\section{Curriculum Development Context}
The legal framework supports education in Slovak Sign Language, and the language was codified in 2020. A news report describes an acute shortage of interpreters and a new interpreting study programme at Trnava University:
\begin{itemize}
\item \url{https://www.noviny.sk/slovensko/557682-tlmocnikov-posunkoveho-jazyka-je-akutny-nedostatok-pomoct-ma-novy-studijny-program}
\end{itemize}
No ready-made K-12 sign-language materials were located, so curriculum development would need to start from the Dictio lexicon, the legal framework and collaboration with deaf organisations and deaf teachers.

\end{document}